\section{Introduction}
\label{sec:intro}

Urban bus systems are the backbone of public transit in many cities. Fleet
operators need near-real-time estimates of which stops and corridors will see
the highest passenger volumes so they can dispatch additional vehicles,
adjust schedules, and allocate drivers. Historically this has been done with
aggregated historical averages and manual reporting, but the availability of
instrumented vehicles (buses with GPS, powertrain sensors, and on-board
passenger counters) has opened the door to data-driven forecasting with
graph and attention models for stop-level demand \cite{li2023pagtsn}.

Modern transit sensor data exhibits a characteristic structural pattern:
dense continuous telemetry (GPS at 1~Hz, speed, power, braking force) arrives
alongside sparse event data (passenger counts updated only when doors open
at stops, typically every 60--120 seconds of driving time). The sparse
series covers approximately 1\% of rows in a raw 1~Hz stream, with the vast
majority of rows carrying only implicit ``no change'' information. This mix
of dense-and-sparse is common in IoT and mobility applications more broadly.

Data scientists approaching this kind of data default to two well-known
preprocessing steps. First, they forward-fill the sparse series onto the
dense timeline, justified by the physical argument that passenger count is
constant between stops (no one boards or alights while the bus is moving).
This transforms 0.9\% usable data into approximately 99\%. Second, they
compute lag features on the now-dense series to capture temporal context,
for example, \texttt{lag(60\,s)} and \texttt{lag(300\,s)} to match typical
rolling-window sizes.

\textbf{An obstacle.} Standard 1~Hz lag features computed on the
forward-filled passenger count series induce target leakage: forward-fill
makes the series piecewise constant between stop events (roughly every
60--120 seconds), and typical lag windows (60~s, 300~s) rarely cross a
stop boundary, so \texttt{lag(60\,s)} equals the current value for
approximately 95\% of rows. Section~\ref{sec:methods-v1} diagnoses the
mechanism in detail; Section~\ref{sec:methods-v2} presents the fix:
compute lag features at stop-event granularity, then forward-fill the
lag values onto the dense grid rather than forward-filling the raw
series first.

\textbf{Baseline.} Our reference point is a standard tree-model pipeline
(Decision Tree, Random Forest, XGBoost) trained on 53 handcrafted sensor
and temporal features (GPS, speed, power, acceleration, time-of-day,
day-of-week, weather) with no passenger count information. On temporal
splits this baseline reaches macro-F1 0.57 for Decision Tree and 0.6066
for Random Forest. The contributions below are measured against these
numbers.

\textbf{Our contributions.}
\begin{enumerate}
  \item We present a strong applied result for stop-level transit demand
    classification on the ZTBus dataset~\cite{widmer2023ztbus}, a
    1,409-mission trolleybus dataset from Zurich, 2019--2022: a Random
    Forest with stop-level lag features reaches temporal macro-F1 =
    0.8287 ($+0.2221$ over the 57-feature cross-sectional baseline) on
    the realistic pre-2022 train, 2022 test evaluation.
  \item We document and fix the forward-fill lag trap underlying this
    result. Adding naively computed 1~Hz passenger-count lag features on
    forward-filled data lifts the Decision Tree from macro-F1 0.57 to
    0.87, but we show this 30-point jump is target leakage: for
    60-second lags on 1~Hz data, approximately 95\% of rows satisfy
    \texttt{lag(60\,s) = current}, so the lag feature is a near-copy of
    the target. Our fix computes lag features at the natural event
    granularity (stop events) and then forward-fills the \emph{lag
    values} onto the dense grid. The resulting features carry genuine
    autoregressive information because they come from strictly earlier
    stops.
  \item We present a negative result on foundation-model feature
    extraction. TimesFM~2.5-200M~\cite{das2024timesfm} penultimate-layer
    embeddings (1280-dim, PCA-reduced to 32) do not improve any model
    beyond the hand-crafted stop-level lags. PCA explained variance
    reaches 100\% at 32 components, indicating the foundation model's
    representations for this domain live in a low-rank subspace. We
    interpret this as the embeddings being largely redundant with
    recent-stop history.
  \item We release a reproducible pipeline implementing all of the
    above, including the event-marker design that enables stop-level
    feature computation through a forward-fill-based preprocessing
    pipeline.
\end{enumerate}

\textbf{Roadmap.} Section~\ref{sec:related} surveys related work.
Section~\ref{sec:data} describes the ZTBus dataset and the classification
task. Section~\ref{sec:eda} presents exploratory analysis.
Section~\ref{sec:methods} develops the methodology, including the leakage
diagnosis and fix. Section~\ref{sec:results} reports quantitative results.
Section~\ref{sec:discussion} discusses generalization and limitations, and
Section~\ref{sec:conclusion} concludes.
